Un resumen es un breve compendio que sintetiza todas las secciones clave de un trabajo de investigación: la introducción, los objetivos, la infraestructura y métodos, los resultados y la conclusión. Su objetivo es ofrecer una visión general del estudio, destacando la novedad o relevancia del mismo, y en algunos casos, plantear preguntas para futuras investigaciones. El resumen debe cubrir todos los aspectos importantes del estudio para que el lector pueda decidir rápidamente si el artículo es de su interés.

En términos simples, el resumen es como el menú de un restaurante que ofrece una descripción general de todos los platos disponibles. Al leerlo, el lector puede hacerse una idea de lo que el trabajo de investigación tiene para ofrecer \cite{elsevier_abstract_2024}.

\textbf{Este informe presenta un análisis experimental sobre el rendimiento práctico de diversos algoritmos y su estrecha relación con la complejidad teórica esperada. El estudio aborda dos problemas computacionales clásicos: el ordenamiento de arreglos unidimensionales (escalando hasta $10^7$ elementos) y la multiplicación de matrices cuadradas (hasta dimensiones de $256 \times 256$). En la primera sección, se evalúan los algoritmos Merge Sort, Patience Sort, Quick Sort y \texttt{std::sort} bajo distintas disposiciones de entrada (aleatoria, ascendente y descendente). En la segunda sección, se contrasta la aproximación iterativa clásica (Naive) con el algoritmo de Strassen, analizando su comportamiento estructural en matrices densas, diagonales y dispersas. A través de la medición exhaustiva de tiempos de ejecución y consumo de memoria, los resultados experimentales no solo validan empíricamente las cotas asintóticas teóricas ($\mathcal{O}$), sino que también evidencian el impacto significativo que tienen las constantes ocultas, la gestión de memoria caché y la naturaleza de los datos sobre la eficiencia algorítmica en la práctica.}
